\chapter{整数}\label{ch:g6:wholes}

数学一切都从计数开始。本章复习整数的书写、比较和运算，
并仔细考察除法。除法这一运算既给出商，
\emph{也}给出
余数。

\section{整数的书写与比较}

\begin{definition}[数位值]\label{def:g6:wholes:place}
我们书写数时使用十个数字（$0$ 到 $9$），并且
一个数字的\emph{数位值}\index{数位值}由它所在的数位决定：在
$5\,208$ 中，数字 $5$ 表示千位，$2$ 表示百位，
$0$ 表示十位（这里没有），$8$ 表示个位：
\[
5\,208 = 5 \times 1000 + 2 \times 100 + 0 \times 10 + 8 .
\]
\end{definition}

\begin{omfigure}
\begin{tikzpicture}[scale=0.9]
  \draw (0,0) grid[xstep=2.3, ystep=0.8] (9.2,1.6);
  \node[font=\small] at (1.15,1.2) {千位};
  \node[font=\small] at (3.45,1.2) {百位};
  \node[font=\small] at (5.75,1.2) {十位};
  \node[font=\small] at (8.05,1.2) {个位};
  \node[font=\large, omDef] at (1.15,0.4) {$5$};
  \node[font=\large, omDef] at (3.45,0.4) {$2$};
  \node[font=\large, omDef] at (5.75,0.4) {$0$};
  \node[font=\large, omDef] at (8.05,0.4) {$8$};
\end{tikzpicture}

{\small $5\,208$ 的数位表。$0$ 十分重要：若没有它，
$5$ 和 $2$ 就会滑到错误的列中，这个数
会被读成 $528$。}
\end{omfigure}

\begin{method}[比较两个整数]\label{met:g6:wholes:compare}
\begin{enumerate}
  \item 位数多的数较大
        ($1\,203 > 989$).
  \item 位数相同时，逐位比较，
        从\emph{左}边开始，第一处不同的数位决定大小：
        $6\,742 > 6\,715$，因为十位上 $4 > 1$。
\end{enumerate}
符号 $<$（“小于”）与 $>$（“大于”）；
符号尖的一端指向较小的数。
\end{method}

\begin{example}\label{ex:g6:wholes:order}
从小到大排列：$908$；$89$；$1\,001$；$980$。
先比较位数：$89$（两位）排在最前，
$1\,001$（四位）排在最后。比较 $908$ 与 $980$：百位都是 $9$，
十位 $0 < 8$，所以 $908 < 980$。最后的答案：
\[
89 < 908 < 980 < 1\,001 .
\]
\end{example}

\begin{omfigure}
\begin{tikzpicture}[scale=1.05]
  \draw[-Stealth] (-0.4,0) -- (10.8,0);
  \foreach \x in {0,1,...,10} \draw (\x,-0.08) -- (\x,0.08);
  \foreach \x/\l in {0/0, 2/200, 4/400, 6/600, 8/800, 10/1000}
    \node[below, font=\small] at (\x,-0.1) {$\l$};
  \fill[omThm] (0.89,0) circle (2pt) node[above, font=\small] {$89$};
  \fill[omDef] (9.08,0) circle (2pt) node[above, font=\small] {$908$};
  \fill[omMeth] (9.8,0) circle (2pt) node[above right=-2pt, font=\small] {$980$};
\end{tikzpicture}

{\small 数排列在一条数轴上：较小的数在左边。每一个刻度
代表 $100$。}
\end{omfigure}

\section{加法、减法、乘法}

\begin{example}[竖式计算]\label{ex:g6:wholes:columns}
做加法或减法时，把个位对齐个位，
十位对齐十位，从右往左计算，需要时进位或退位：
\[
\begin{array}{r}
4\,5\,7 \\
+\ 2\,8\,6 \\
\hline
7\,4\,3
\end{array}
\qquad\qquad
\begin{array}{r}
6\,0\,3 \\
-\ 1\,4\,7 \\
\hline
4\,5\,6
\end{array}
\]
做加法时：$7 + 6 = 13$，写 $3$，进 $1$；然后
$5 + 8 + 1 = 14$，写 $4$，进 $1$；再算 $4 + 2 + 1 = 7$。
用加法检验减法：$456 + 147 = 603$。
\end{example}

\begin{remark}[词汇]
加法的结果是\emph{和}\index{和}，减法的结果是
\emph{差}\index{差}，乘法的结果是
\emph{积}\index{积}，除法的结果是
\emph{商}\index{商}。参与运算的数称为
\emph{项}（对于 $+$、$-$）或\emph{因数}（对于 $\times$）。
\end{remark}

\begin{example}[巧算]\label{ex:g6:wholes:smart}
调换项或因数的顺序，常常能省去不少计算：
\[
17 + 58 + 3 = (17 + 3) + 58 = 20 + 58 = 78 ,
\]
\[
25 \times 17 \times 4 = (25 \times 4) \times 17 = 100 \times 17 = 1700 .
\]
\end{example}

\section{除法}

\begin{theorem}[带余除法]\label{thm:g6:wholes:division}
给定两个整数 $a$（\emph{被除数}）和 $b \neq 0$（
\emph{除数}），有且仅有一种方式将其写成
\[
a = b \times q + r
\qquad\text{其中}\qquad
0 \leq r < b .
\]
$q$ 称为\emph{商}\index{商}，$r$ 称为
\emph{余数}\index{余数}，它们是 $a$ 除以 $b$ 所得的结果。
\end{theorem}
\admitted

\begin{example}\label{ex:g6:wholes:division}
逐步计算 $137$ 除以 $4$。
\begin{enumerate}
  \item $4$ 在 $13$ 里包含几次？答：3 次
        （$4 \times 3 = 12$），余数 $13 - 12 = 1$。
  \item 落下 $7$，$4$ 在 $17$ 里包含几次？
        答：4 次（$4 \times 4 = 16$），余数 $1$。
\end{enumerate}
于是 $q = 34$，$r = 1$：
\[
137 = 4 \times 34 + 1, \qquad\text{且显然 } 0 \leq 1 < 4 .
\]
验算：$4 \times 34 = 136$，且 $136 + 1 = 137$。
\end{example}

\begin{omfigure}
\begin{tikzpicture}[scale=0.95]
  \draw[-Stealth] (-0.3,0) -- (11.2,0);
  \foreach \i in {0,1,...,34} \draw ({\i*0.32},-0.06) -- ({\i*0.32},0.06);
  \foreach \i in {0,10,20,30}
    \node[below, font=\footnotesize] at ({\i*0.32},-0.08) {$4\times\i$};
  \draw[omDef, very thick] (0,0.14) -- (10.88,0.14);
  \draw[omThm, very thick] (10.88,0.14) -- (10.96,0.14);
  \fill (10.96,0) circle (1.8pt) node[above, font=\small] {$137$};
  \node[omDef, font=\small] at (5.4,0.55) {$34$ 组，每组 $4$ 个};
\end{tikzpicture}

{\small 计算 $137$ 除以 $4$：可以装满 $34$ 组，每组 $4$ 个
（共 $136$），还剩 $1$ --- 余数总是
比除数小。}
\end{omfigure}

\begin{definition}[能被……整除]\label{def:g6:wholes:divisible}
当余数为 $0$ 时，称 $a$ 能被 $b$
\emph{整除}\index{整除}：例如 $135$ 能被 $5$ 整除，
因为 $135 = 5 \times 27 + 0$。
\end{definition}

\begin{method}[含除法的应用题]\label{met:g6:wholes:problems}
算出 $a = bq + r$ 后，务必回到题目所问：
\begin{enumerate}
  \item “能装满几箱/几队/几页？” --- 答案是
        商 $q$；
  \item “还剩几个？” --- 答案是
        余数 $r$；
  \item “所有人需要几箱？” --- 若 $r > 0$，则需要
        比商\emph{多}一箱，即 $q + 1$。
\end{enumerate}
\end{method}

\begin{example}\label{ex:g6:wholes:bus}
$137$ 名学生乘坐每车 $40$ 个座位的客车外出，列除法算式：
$137 = 40 \times 3 + 17$。3 辆客车可载 $120$ 名学生，剩下 $17$
名学生 --- 他们也需要一辆车！所以需要 $4$ 辆客车，
尽管商是 $3$。
\end{example}

\section{练习}

\begin{exercise}[$\star$]\label{exo:g6:wholes:1}
用数字写出：“三千零四十七”；“二万
零五百”；“一百零四”。再用汉字写出 $60\,013$。
\end{exercise}

\begin{exercise}[$\star$]\label{exo:g6:wholes:2}
在数 $47\,209$ 中：百位上的数字是几？个位
上的数字是几？数字 $4$ 表示多少个什么？
把该数写成 $1$、$10$、$100$、\dots\ 的倍数之和，仿照
\cref{def:g6:wholes:place}。
\end{exercise}

\begin{exercise}[$\star$]\label{exo:g6:wholes:3}
抄写下列各式，并在空处填入 $<$ 或 $>$：
\[
507 \;?\; 570, \qquad
1\,099 \;?\; 989, \qquad
23\,456 \;?\; 23\,465, \qquad
9\,999 \;?\; 10\,000 .
\]
\end{exercise}

\begin{exercise}[$\star$]\label{exo:g6:wholes:4}
将下列各数从小到大排列：$2\,043$；$978$；$2\,304$；$2\,034$；
$1\,001$。
\end{exercise}

\begin{exercise}[$\star$]\label{exo:g6:wholes:5}
用竖式计算：$348 + 576$；\quad $805 - 268$；\quad
$307 \times 26$。用加法检验减法的结果。
\end{exercise}

\begin{exercise}[$\star$]\label{exo:g6:wholes:6}
运用巧算，将加数或因数适当分组后计算：
\[
38 + 45 + 12 + 5, \qquad
2 \times 83 \times 50, \qquad
4 \times 78 \times 25 .
\]
\end{exercise}

\begin{exercise}[$\star$]\label{exo:g6:wholes:7}
对每个除法，写出商与余数，
并写出等式 $a = b \times q + r$：
\[
95 \div 7, \qquad
230 \div 6, \qquad
504 \div 8 .
\]
\end{exercise}

\begin{exercise}[$\star$]\label{exo:g6:wholes:8}
不做除法，说明等式 $173 = 8 \times 20 + 13$ 为何
\emph{不是} $173$ 除以 $8$ 的欧几里得除法，并求出正确的
商与余数。
\end{exercise}

\begin{exercise}[$\star\star$]\label{exo:g6:wholes:9}
鸡蛋按每盒 $12$ 个装箱。
\begin{enumerate}
  \item 某农场有 $2\,000$ 个鸡蛋。可以装满多少盒，
        还剩多少个鸡蛋？
  \item 某面包店需要 $150$ 个鸡蛋，需要买多少盒？
\end{enumerate}
\end{exercise}

\begin{exercise}[$\star\star$]\label{exo:g6:wholes:10}
今天是星期二。再过 $100$ 天是星期几？
（一周有 $7$ 天：用除法，并考虑余数。）
\end{exercise}

\begin{exercise}[$\star\star$]\label{exo:g6:wholes:11}
在一个除以 $9$ 的除法中，商是 $56$，余数
为可能取到的最大值。被除数是多少？
\end{exercise}

\begin{exercise}[$\star\star\star$]\label{exo:g6:wholes:12}
我是一个三位数。我的百位数字是个位数字的两倍，
十位数字是 $0$，各位数字之和是 $9$。我是哪个数？
（逐步推理，然后检验。）
\end{exercise}

\section{问题：小高斯与前一百个
数的和}

\begin{problem}[{周末问题 --- 十秒内计算 $1 + 2 + 3 + \dots +
100$ 的和，以及三角形数}]
\label{pb:g6:wholes:1}
一个著名的故事：为了让全班学生忙碌起来，一位校长让学生
把 $1$ 到 $100$ 的所有数加起来。一个男孩
几乎立刻在石板上写下 $5\,050$ --- 卡尔·弗里德里希·高斯，九岁，
后来成长为史上最伟大的数学家之一。
他的秘诀无非是\emph{巧算}
（\cref{ex:g6:wholes:smart}）：把各项分组，使每一组
都易于计算。本题重新发现了他的方法，并使其适用于每一种
情形，同时说明这类和在日常生活中出现得多么频繁。

\medskip
\noindent\textbf{第一部分 --- 高斯的配对。}
\begin{enumerate}
  \item 通过巧妙分组，计算 $1 + 2 + 3 + 4 + 5$。
  \item 对于 $1 + 2 + 3 + \dots + 10$，将第一项与最后一项配对
        （$1 + 10$），第二项与倒数第二项配对
        （$2 + 9$），依此类推。每一对的值是多少？共有
        多少对？算出最后的结果。
  \item 对 $1 + 2 + \dots + 20$ 用同样的方法：每一对的值是
        多少，共有多少对，总和是多少？
  \item 现在计算校长布置的和 $1 + 2 + \dots + 100$：说明
        为什么配对 $1 + 100$、$2 + 99$、$3 + 98, \dots$ 的值都相同，
        并说明为什么恰好有 $50$ 对，
        最后一对是 $50 + 51$。
  \item 将高斯的答案写成一个乘法算式，
        并计算结果。
\end{enumerate}

\medskip
\noindent\textbf{第二部分 --- 永不失效的技巧。}
\begin{enumerate}[resume]
  \item 对 $1 + 2 + \dots + 9$ 试用配对法：这次有一个
        数没有配对的伙伴。是哪一个？以它为中心配对
        并计算和。
  \item 下面是该技巧的一个版本，无论项数为偶数还是奇数都适用：把这个和写两次，
        一次正着写，一次倒着写，
        一行放在另一行下面，
        \[
        \begin{array}{ccccccc}
        1 & + & 2 & + & \dots & + & 9 \\
        9 & + & 8 & + & \dots & + & 1
        \end{array}
        \]
        然后\emph{逐列}相加。每一列的值是多少？
        一共有多少列？说明为什么两行之和为
        $9 \times 10 = 90$，由此可得
        $1 + 2 + \dots + 9 = 45$ --- 与
        第 6 题的答案相同。
  \item 用两行技巧计算 $1 + 2 + \dots + 50$，
        再计算 $1 + 2 + \dots + 200$。
  \item 计算前一百个\emph{偶数}的和，
        即 $2 + 4 + 6 + \dots + 200$。（将每一项与高斯和中的
        对应项比较，无需新的配对。）
  \item 计算 $101 + 102 + \dots + 200$，利用你已经
        会算的两个和以及一次减法。
\end{enumerate}

\medskip
\noindent\textbf{第三部分 --- 处处可见的三角形数。}
$1$，$3$，$6$，$10$，$15, \dots$ 这些数，即和
$1 + 2 + \dots + n$ 的结果，被称为\emph{三角形
数}\index{三角形数}，因为它们计数的是
按三角形堆叠的物体。
\begin{enumerate}[resume]
  \item 八位朋友相聚，每人与其他每位
        恰好握手一次。说明握手次数为何是
        $7 + 6 + 5 + 4 + 3 + 2 + 1$，并计算结果。
  \item 杂货商用罐头堆一个三角形棱锥：底层 $12$ 罐，
        第二层 $11$ 罐，依此类推，直到顶部
        只有一罐。她需要多少罐？（使用
        第 7 题的技巧。）
  \item 将和 $1 + 2 + 3 + 4$ 画成一个正方形阶梯：
        第一列一个正方形，第二列两个，第三列三个，
        第四列四个。在你的图上说明
        \emph{两个}这样的阶梯，其中一个上下颠倒，
        恰好拼成一个 $4$ 行、每行
        $5$ 个正方形的长方形 --- 并解释为何这正是第 7 题的两行
        技巧，只是改用图形表示。
  \item 钟表在一点整敲 $1$ 下，两点整敲 $2$ 下，
        \dots，十二点整敲 $12$ 下。十二小时共敲
        多少下？那么整整一天
        呢？
  \item 最后一题献给高斯：$5\,050$ 秒 --- 比一个半小时多还是
        少？将它换算为小时、
        分钟和秒，使用欧几里得除法
        （\cref{thm:g6:wholes:division}）两次。
\end{enumerate}
\end{problem}
